\chapter{Önsöz}
``Peçete'' ismi benim aşağıdaki sözümden geliyor.
\begin{quote}
\small
Hâlâ liseye giden birkaç arkadaşla kısa bir yemek yiyorum, ve 
onlar son haftalarda nelerle uğraştığımı soruyorlar. Ben kategori teorisi
öğrendiğimi açıklıyorum, onlar kategori teorisinin ne olduğunu soruyor.
Onlara, kategori teorisinin objeleri sadece yapıyı koruyan fonksiyonlara bakarak 
soyutlamakla alakalı olduğunu söylüyorum, ve sonra standart $\catname{Grp}$ 
örneğini vermeye çalışıyorum, ama homomorfizmanın ne olduğunu bilmedikleri 
ortaya çıkıyor. Homomorfizmanın ne olduğunu anlatmaya başlıyorum, 
grubun ne olduğunu bilmedikleri ortaya çıkıyor. Sonra grubun ne olduğunu anlatmaya 
başlıyorum, ama grup aksiyomlarını peçeteme yazmayı bitirene kadar neden gruplardan 
bahsettiğimi unutuyorlar. Sonra saat bir oluyor, insanların gideceği yerler var malum, ve 
ben sadece şunu düşünüyorum:\\[0.5ex]
\emph{``Aslında kırk dakika değil de kırk saatim olsaydı bunların 
hepsini anlatabilirdim''.}
\end{quote}
Başlarda bu kitap o kırk saat olacaktı, fakat o günden bu yana
çok daha büyüdü.

\section*{Kitap hakkında}
\emph{Sonsuz bir Peçete}, ileri matematiğin büyük bir kısmına 
hafif ama çoğunlukla bağımsız bir giriştir.

Hemen söylemeliyim ki, bu kitap konularda özel
kitap ve kursların yerini alamaz, derinlikleri kıyaslanamaz.
Ancak bu ``hafif'' yaklaşımın iyi bir yanı ise materyal daha büyük 
bir kitle için erişilebilirlik kazanıyor, amaç tam bir sömestir 
derslerinin yerini almaktansa okura herhangi bir konuya 
geniş bir bakış vermek olduğu için.

İlk başta bu kitabı yetenekli liselileri düşünerek---özellikle 
olimpiyat geçmişi olanlar için---yazmıştım. Bu kültürel yanlılık 
kitabın çok yerinde hissedilebiliyor, ki özellikle matematik 
yarışmalarından alınan çeşitli problemlerde. 
Fakat bence bu kitap genelinde meraklı ve matematiksel olgunluğu olan 
herhangi bir kimsenin kullanabileceği bir referans.
Örneğin (kesinlikle bunlara sınırlı olmayacak şekilde): 
lisans düzeyinde matematik okuyanlar, fizik veya bilgisayar bilimi 
okuyanlar, kendi alanı dışındaki alanları merak eden doktora öğrencileri, 
matematiği seven ama matematik yarışmalarını sevmeyen ilerlemiş liseliler, 
ve Türkçe bilen zeki kediler. 

\section*{Source code}
Bu proje GitHub'da
\url{https://github.com/vEnhance/napkin} adresinde bulunmakta.
Pull request atmaktan çekinmeyin! 
Öneriler ve hatalar için ayrıca e-posta de gönderebilirsiniz.
(Dil hataları için lütfen Evan değil benimle iletişime geçin).

\section*{Peçete yaklaşımının felsefesi}
Benim bildiğim kadarıyla liseliler için ileri matematiğin iki
şekli vardır:
\begin{itemize}
    \ii Biri sana der ki, ``küresel bir kedinin saçını tarayamazsın'' 
        ama neden doğru olması gerektiğini, ``saçını tarama''nın aslında 
        ne anlama geldiğini, altında yatan teoriyi söylemez, senin aklında 
        sadece kafanda yarım yamalak bir kavram olarak bırakır.
    \ii Bir ders alırsın ve her sonucu tek tek her detayıyla ispatlarsın, 
        ve bir yerden sonra hocanın dediği umrunda olmaktan çıkar.
\end{itemize}

İlk yaklaşımın ne kadar rahatsız edici olduğunu biliyorsunuzdur, o yüzden 
de ikinci yaklaşım yaygın olarak kullanılıyor. Ama bence bu ikisi arasında bir orta
nokta bulunması gerekiyor.

%I was talking to a friend of mine one day who described briefly
%what the Israel IMO training looked like.
%It turns out that rather than actually preparing for the IMO,
%the students would, say, get taught a semester's worth of
%undergraduate algebra in the couple weeks.
%This might seem absurd, but I think if your goal is to just show the students
%what this algebra thing they keeping hearing about is,
%and your students have substantially more mathematical maturity
%relative to their knowledge (e.g.\ are IMO medalists),
%it seems like you really could make a lot of headway.

%For example: often classes like to prove things for completeness.
%I personally find that many proofs don't really teach anything,
%and that it is often better to say
%``you could work this out if you wanted to,
%but it's not worth your time''.
Bu kitabın amacı üniversitedeki gibi egzersiz çözüp problem çözmeyi 
öğretmek değil,\footnote{Bu demek değil ki problem çözme 
    değerli değil; kalben olimpiyat koçuyum, 
    ama bu kitabın amacı değil.}
alanda araştırmaya hazırlamak da değil. Ben sadece 
matematikte enteresan şeyler göstermek istiyorum. 
Gösterilen şeyler hatırlanabilir ve önemli şeyler 
olmalı. 
Bu yüzden çoğu ders kitabında bütünlük adına 
yapılan ispatlar bu kitapta yok---sadece görmeye değer 
olduğunu düşündüğüm bir fikir varsa gösteririm. 
Bilhassa bir teoremin neden doğru \emph{olması gerektiğini}
anlatmaya, onun ispatını vermeye göre daha çok vurguluyorum.
Bu kitapta çok rastlanılan bir tema:
\begin{moral}
    Doğal açıklamalar ispatlardan daha değerlidir.
\end{moral}

Umudum birinin Peçete'den herhangi bir bölüm okuduktan sonra 
şunları elde etmesi:
\begin{itemize}
    \ii Ana karakterlerin tam tanımlarının bilgisi,

    \ii En önemli örneklerle haşır neşir olma,

    \ii Meşhur teoremlerin tam ne dediklerininin bilgisi 
    ve neden doğru olduğunu bilme hissi
\end{itemize}
Bir şeyin ``neden'' doğru olduğunu bilmenin çok şekli vardır.
Bazen tam, kesin bir ispat gerektirir, bazen ispatın 
sadece fikri verilir. Bazen ise sadece birkaç 
önemli örnek ve biraz tezahürat yeter.

Tabii ki bu ilgili alanda araştırma yapmak istiyorsan yeterliye 
yakın bile değil, ama uzman olmayan bir meraklı için 
umarım Vikipedi'den ve kısa anlatılardan daha etkili olur. 
Bir konuyu ciddi derinlikte öğrenmek istiyorsan bile
umarım ki bu kitap tam dalmadan önce iyi bir kuşbakışı 
perspektif verir. Çünkü aslında çoğu yönden bu materyal 
``keşke başlamadan bana anlatılsaydı'' dediğim şeylerden 
ibaret.

%In terms of learning math, I suspect that
%there is some sense in which reading a textbook
%is like watching a movie.
%The first time you watch it, you spend most of your bandwidth
%getting to know the characters,
%keeping up with what's going on,
%trying to predict where it's headed next.
%But if you like the movie enough and watch it a second
%or third, or $n$th time, you'll pick up a lot of things
%that you didn't notice the first time through
%(e.g.\ all the hints that Z was the main villain all along).

\section*{Daha çok pedagojik yorum ve referans}
Eğer her bölüm için tek tek pedagojik seçimlerimi
açıklasaydım önsöz çok uzun olurdu, o yüzden bu 
yorumların çoğu \cref{ch\string:refs}'da.

Özellikle sıklıkla belirli referanslar kullanırım, ve 
bu ekin sonunda daha çok referans var. Konularda
daha da ilerlemek istiyorsanız buraya bakmak iyi olabilir.

\section*{Tarihi ve kişisel notlar}
Bu kitabı 2014'ün Aralık'ında, Harvard'da ilk lisans 
sömestirimi bitirdikten sonra yazmaya başladım. 
İleri matematiğe girdikçe bundan sonraki 18 
ay boyunca ana uğraşım oldu.
Neredeyse sadece matematik dersleri aldım (diğer tüm 
zorunlu dersleri görmezden gerek), ve bunları toplayabildiğim kadar 
Peçete'de topladım (kendi zamanımda öğrendiğim matematiği de ekledim). 

2016'nın Ağustos'una doğru sonunda yoruldum. 
O aralar ilk taslaklar internete yayınlanmıştı, ve ben bir geri adım attım.
Biraz yorulmuştum, ve lisans eğitimimin kalan iki yılını
Amerika matematik olimpiyat takımının koçu olarak ve 
herkes mezun olup ayrı yollarına gitmeden arkadaşlarımda 
olabildiğince sık görüşerek geçirdim\footnote{Alternatif: ``\dots Kalan 
    iki yılı öğrendiğim her şeyi unutarak geçirdim.'' İyi ki 
    Peçete'deki notları yazmışım dedirtti.}.

% It was also at this time I started to get into Korean pop dance,
% which was the first (and only) extracurricular activity
% I did during my entire undergraduate stay.

Bu iki yıl süresince okurlar bana çok sayıda güzel söz, teşekkür, hata tahsisleri ve 
öneri gönderdi. 2018'in Kasım'ında, doktora öğrenciliğimin ilk sömestirinde, 
başladığımı bitirmeye karar verdim. Birkaç\todo{whether to add joke about 
this or stick to the script} ay sonra elimdeki bu.

\section*{Teşekkürler}
Bu kitap için sayısız kişiye müteşekkirim. Bı (kesinlikle eksik) kısmi 
bir liste.


\begin{itemize}
\ii Tüm öğretmen ve hocalarıma bu notlarda yer alan çoğu materyali 
öğrettikleri için, ve kitapta alıntı yaptığım her yazara teşekkür ederim. 
Özellikle \cite{ref:55a}, \cite{ref:msci},
\cite{ref:manifolds}, \cite{ref:gathmann}, \cite{ref:18-435},
\cite{ref:etingof}, \cite{ref:145a}, \cite{ref:vakil},
\cite{ref:pugh}, \cite{ref:gorin} etkililerdi.

\ii Önizleme taslaklarını okuyan, hata bulan, öneride bulunan
onlarca arkadaş ve yabancıya teşekkür ederim. Özellikle 
binden fazla hata bulan Andrej Vukovi\'c ve Alexander Chua'ya 
teşekkür ederim. Ayrıca çok tahsiste bulunan Brian Gu ve Tom Tseng'e 
teşekkür ederim (Hata bulduysanız veya öneriniz varsa çok duymak isterim).
Royce Yao ve \texttt{user202729}'a yazdıkları misafir bölümler için
teşekkür ederim.

\ii Ön kapağı çizen Jenny Chu ve Lanie Deng'e teşekkür ederim.  

\ii Bu projenin gelişim aşamasında aldığım çok sayıda güzel söz için 
şükranımı belirtmek isterim. Bu cömert yorumlar kitabın üzerine 
çalışmamı sürdürdü ve çoğunlukla onlar sayesinde 2018'in Kasım'ında 
Peçete'yi yeniden güncellemeye başladım.
\end{itemize}

Son olarak, matematik olimpiyat topluluğuna 
(benim ve Peçete'nin anayurdu) kocaman bir teşekkür iletiyorum.
Gördüğüm onca heves, teşvik ve teşekkür sayesinde bu kitabı yazmaya karar verdim. 
Öbür türlü hiç bir şekilde böyle bir projenin yapılabileceğini düşünecek 
kadar ukala olamazdım. Ayrıca kesinlikle şu an olduğum yerde,
IMO hayallerimin peşine düşerken 
geçirdiğim hayat değiştiren serüven olmasa olamazdım. 
Sonsuza dek TWN2!

